\section{Introduction}
Power amplifiers (PAs) are a critical subsystem in any RF transmitter chain, responsible for boosting the signal to the power levels required for reliable communications over long distances. In satellite systems, this role is particularly demanding, as the PA must deliver sufficient power to close the downlink budget while operating under strict energy, mass, and thermal constraints, and withstanding the harsh conditions of the space environment, including ionizing radiation, wide thermal swings, and high vacuum.

Gallium Nitride High Electron Mobility Transistors (GaN HEMTs) have established themselves as strong competitors to conventional technologies such as GaAs and Si in RF power amplifiers designs, given their superior breakdown voltage, elevated power density, wide thermal operating range, and capability for operation at high frequencies \cite{Boomer2020, Mishra2008, Deblecker2015}. These properties make GaN transistors well suited for a wide range of applications, such as broadband, high-power, and high-efficiency designs.

In parallel, the CubeSat standard has fundamentally changed the landscape of space missions. Originally conceived as a low-cost platform based on 10×10×10\,cm modular units (1\,U) \cite{cubesat_design_spec}, CubeSats have evolved into capable research and commercial platforms used for several applications, such as Earth observation, telecommunications, and technology demonstration. Their reduced size and mass impose severe constraints on the RF subsystems: the PA must be compact, efficient, and capable of sustaining reliable operation in LEO without active thermal management or human intervention.

This work presents the design and experimental characterization of a single-stage, Class AB, 1\,W GaN power amplifier targeting the S-Band downlink at 2.2\,GHz for CubeSat applications. The output power level was determined from a link budget analysis supporting data rates up to 4.44\,Mbps.

Building on well-established techniques (Class AB biasing, load-pull-based matching, resistive stabilization), this work proposes a link-budget-driven, low-cost FR4 design tailored to CubeSat downlinks, validated through full experimental characterization and an analysis of the simulation-to-measurement frequency shift.

This paper is organized as follows. Section~II presents the link budget analysis and theoretical background. Section~III describes the proposed circuit design, Section~IV the measurement setups, and Section~V the experimental results. Section~VI draws conclusions.